\section{Results}
\label{sec:results}

\pgfplotstableset{col sep=comma}
\subsection{Experimental setup}
\label{sec:setup}
\textbf{Data and labels.} I use the per-node release of M100 ExaData \cite{borghesi2023m100}: one Parquet file per node with 15-minute windows, the average, standard deviation, minimum and maximum of each IPMI metric, and a Nagios label from the scheduler-availability check, which reports whether the scheduler can place jobs on the node (\cref{tab:nagios}). The check runs every 15 minutes and always records a state: 87.2\% of the 12.35\,M node-windows in the evaluation set carry an explicit OK, 2.0\% critical and 0.04\% unknown. The remaining 10.7\% have no label. They occur in whole blocks (March--April 2020 entirely, 91\% of March 2021) and are excluded from every metric.

\begin{table}[t]
\caption{Nagios state of the scheduler-availability check.}
\label{tab:nagios}
\centering\small
\begin{tabular}{cp{6.4cm}}
\toprule
State & Meaning \\
\midrule
0 & OK: the node can accept jobs \\
2 & Critical: drained or unavailable to the scheduler (the anomaly label) \\
3 & Unknown: the check could not determine the state; 39 timestamps, only during cluster-wide events, counted as anomalous \\
missing & Unmonitored: no label, excluded \\
\bottomrule
\end{tabular}
\end{table}

\textbf{Slices.} \Cref{tab:setup} lists the data used. The streaming replay covers four weeks of the three downloaded racks. For evaluation over a long history I read only the label columns of all 49 racks (857 nodes) with HTTP range requests and kept the 8 racks with the most independent events, 146 nodes in total. These racks were selected on label counts over the whole history, so event rates are higher than for a random rack.

\begin{table}[t]
\caption{Data slices. Windows are monitored node-windows.}
\label{tab:setup}
\centering\small
\begin{tabular}{p{1.5cm}p{6.3cm}}
\toprule
Use & Content \\
\midrule
Stream replay & 49 nodes (racks 10--12), 2020-06-14 to 2020-07-12, 6.32\,M events (48 values per node and window) \\
Detection & 146 nodes. Train 2020-03 to 08, validation (thresholds) 2020-09 to 12, test 2021-01 to 2022-09: 8.01\,M windows, 2.16\% drained, 517 independent events \\
Anticipation & 146 nodes, chronological 80/20 split, test 2022-03-26 to 2022-09-28: 1,313 onsets at 1\,h among 2.54\,M rows (base rate 0.05\%), 1,242 at 6\,h \\
\bottomrule
\end{tabular}
\end{table}

\textbf{Events and metrics.} An \emph{event} is a run of at least four consecutive drained windows on a node preceded by a normal window; onsets at the same timestamp within one rack count once. An onset is \emph{cluster-wide} when at least half of the eligible nodes have an onset at the same window, and \emph{node-specific} otherwise; 70\% of the 1\,h test onsets are cluster-wide. I report ROC-AUC pooled over nodes and as the per-node median or mean, PR-AUC, precision and F1 at the threshold with the best validation F1, and for anticipation the number of false alerts per node-day at a fixed recall. Intervals are bootstrap 95\% intervals over nodes. Models are compared with \emph{paired} per-node AUC differences, which cancel node-to-node variation.

\textbf{Single-feature rule.} As a training-free baseline I score each window by one raw feature. The feature and its sign are chosen \emph{on training data only}, among all 48 features, by mean per-node AUC; this picks $-$\texttt{p0\_mem\_power\_std} for detection, $-$\texttt{total\_power\_max} for anticipation at 1\,h and $-$\texttt{gpu0\_core\_temp\_avg} at 6\,h. A rule that I first hand-picked ($-$total power) after inspecting all data scores 0.728 and 0.654 at 1\,h and 6\,h, slightly better than the selected one; it is shown separately as an optimistic reference.

\subsection{Pipeline (RQ1)}
\label{sec:pipeline}
\Cref{tab:pipeline} reports three replays of the event stream on one machine (single broker, Spark in local mode, 5-second micro-batches), each on a fresh stack. The producer sustains 74\,k messages/s. Spark keeps up at every rate: a micro-batch of $\sim$380\,k rows takes 1.9\,s, results land in TimescaleDB 3.3\,s (median) and 5.9\,s (p95) after publish, and the job finishes within 3--5\,s of the producer. Most of that latency is the 5-second trigger interval.

\begin{table*}[t]
\caption{Throughput and latency at three replay rates (fresh stack per run). Detector lag is scoring time minus publish time.}
\label{tab:pipeline}
\centering\small
\begin{tabular}{lrrrr}
\toprule
Replay & msg/s & Batch (s) & End-to-end p50/p95 (s) & Detector lag p50 (s) \\
\midrule
Max speed, 6.3M events & 74k & 1.9 & 3.3 / 5.9 & 75.6 \\
20 windows/s, 1.6M & 47k & 1.7 & 3.7 / 6.0 & 14.2 \\
5 windows/s, 1.6M & 12k & 0.9 & 3.2 / 5.4 & 0.8 \\
\bottomrule
\end{tabular}

\end{table*}

The Python detector service is the bottleneck: its lag is 0.8\,s at 12\,k msg/s but grows to 14\,s at 47\,k and 76\,s at 74\,k msg/s, identically for a half-space-tree and a label-based detector, so the cost is decoding and vector reassembly, not scoring. The topic has eight partitions, so several detector instances could share the nodes.

\emph{Restart.} I killed Spark about 75 s into a full replay and restarted it 15 s later. It resumed from its checkpoint, worked off the backlog in 500\,k-row batches and, compared with an independent recomputation, produced exactly the expected 1,580,544 hourly windows with no missing, extra or duplicated key and no differing count, average, minimum or maximum.

\subsection{Detecting drained nodes (RQ2)}
\label{sec:detection}
\Cref{tab:detection} compares generic unsupervised detectors with the references on the test period. All six unsupervised detectors are poor: per-node median AUC at most 0.59, pooled AUC between 0.34 and 0.51, F1 at most 0.055. Relative entropy and Pearson, with Rettig et al.'s thresholds, flag 90--98\% of all windows. Half-space trees from CapyMOA are the best, and the peer-relative variants do not help.

\begin{table*}[t]
\caption{Detection of the current state, test 2021-01 to 2022-09. AUC is pooled, per-node AUC is the median over 146 nodes; F1 uses the best validation threshold.}
\label{tab:detection}
\centering\small
\begin{tabular}{lrrrr}
\toprule
Method & AUC & Per-node AUC & PR-AUC & F1 \\
\midrule
Always OK (flag nothing) & 0.500 & 0.500 & 0.022 & 0.000 \\
Best single feature, chosen on training data (-p0\_mem\_power\_std) & 0.818 & 0.819 & 0.089 & 0.175 \\
Supervised GBT (ceiling) & 0.881 & 0.859 & 0.372 & 0.362 \\
\midrule
HST (CapyMOA) & 0.447 & 0.592 & 0.024 & 0.055 \\
HST (River) & 0.338 & 0.454 & 0.017 & 0.041 \\
HST, peer-relative (CapyMOA) & 0.393 & 0.457 & 0.018 & 0.040 \\
HST, peer-relative (River) & 0.363 & 0.426 & 0.019 & 0.038 \\
Relative entropy & 0.448 & 0.503 & 0.019 & 0.042 \\
Pearson vs peers & 0.514 & 0.455 & 0.026 & 0.041 \\
\bottomrule
\end{tabular}

\end{table*}

Every unsupervised detector is clearly worse than the single-feature rule in the paired test, by 0.25 to 0.39 per-node AUC (winning on 0--3\% of nodes). The supervised GBT is better than the rule by +0.047 [+0.039, +0.055], winning on 90\% of nodes. This holds against the stronger hand-picked power rule as well (+0.04): the GBT wins on all 8 racks and in 15--16 of 20 test months, and three training seeds differ by 0.002.

\begin{figure}[t]
\centering
\begin{tikzpicture}
\begin{axis}[width=\columnwidth, height=4.4cm, xlabel={total power / node idle level}, ylabel={share of windows},
  xmin=0.8, xmax=3, ymin=0, ymax=0.6, legend style={font=\scriptsize, at={(0.97,0.95)}, anchor=north east},
  tick label style={font=\scriptsize}, label style={font=\footnotesize}, every axis plot/.append style={thick}]
\addplot[cmodel, const plot, mark=none] table[col sep=comma, x=x, y=normal] {results/power_hist.csv};
\addplot[crule, const plot, mark=none] table[col sep=comma, x=x, y=drained] {results/power_hist.csv};
\legend{normal, drained}
\draw[dashed, gray] (axis cs:1.1,0) -- (axis cs:1.1,0.6);
\end{axis}
\end{tikzpicture}
\caption{Power relative to each node's idle level (10th percentile of training windows) in normal and drained test windows; the last bin collects values above 3. The dashed line marks the idle band (110\%).}
\label{fig:power}
\end{figure}

\emph{Why.} \Cref{fig:power} shows that 73\% of drained windows sit at idle power, against 6\% of normal windows. The converse is weak: of all idle windows only 21\% are drained, about ten times the base rate but far from decisive, and 79\% are nodes that are available and simply have no jobs. Power ranks windows well but cannot classify them: even the supervised GBT reaches only 33\% precision at 39\% recall. Density-based detectors rank \emph{busy} windows, which are heterogeneous, as more unusual than the dense idle mode where drained nodes sit.

\emph{RUAD replication.} A per-node LSTM autoencoder on all 352 aggregates of 32 random nodes scores a pooled AUC of 0.569 (per-node mean 0.606), which reproduces the 0.57 that the dataset authors report for this public release and is far below the 0.767 of the original paper, which uses non-public features and cleaned labels. On the same nodes and windows, CapyMOA half-space trees score 0.726, but across all 146 nodes they score 0.45 pooled (0.59 per-node median), so a sample of 32 nodes is misleading.

\subsection{Anticipating drains (RQ3)}
\label{sec:anticipation}
\Cref{tab:anticipation} gives the main results at 1\,h and 6\,h with the paired difference to the training-selected rule.

\begin{table*}[t]
\caption{Anticipation of drains within 1\,h and 6\,h. AUC is the mean per-node AUC with bootstrap 95\% interval; $\Delta$ is the paired difference to the best single-feature rule$^a$ (all onsets, 1\,h); alerts are false alerts per node-day at 25\% recall. $^a$ chosen on training data ($-$total\_power\_max at 1\,h, $-$gpu0\_core\_temp\_avg at 6\,h). $^b$ chosen after inspecting all data, optimistic.}
\label{tab:anticipation}
\centering\small
\begin{tabular}{lrrrrrr}
\toprule
Model & 1\,h AUC & 1\,h node-spec. & $\Delta$ vs rule & 1\,h alerts & 6\,h AUC & 6\,h alerts \\
\midrule
Always OK & 0.500\,{\scriptsize[.500,.500]} & 0.500\,{\scriptsize[.500,.500]} & -- & 96.0 & 0.500\,{\scriptsize[.500,.500]} & 96.0 \\
Best single feature (train)$^a$ & 0.721\,{\scriptsize[.707,.735]} & 0.731\,{\scriptsize[.690,.769]} & -- & 4.9 & 0.637\,{\scriptsize[.621,.652]} & 11.2 \\
Low power, hand-picked$^b$ & 0.728\,{\scriptsize[.713,.741]} & 0.736\,{\scriptsize[.697,.772]} & -- & 3.6 & 0.654\,{\scriptsize[.638,.669]} & 6.5 \\
GBT on power only & 0.695\,{\scriptsize[.680,.708]} & 0.689\,{\scriptsize[.653,.726]} & -- & 5.1 & 0.621\,{\scriptsize[.607,.635]} & 7.5 \\
\midrule
GBT, telemetry + peers (offline) & 0.739\,{\scriptsize[.726,.753]} & 0.710\,{\scriptsize[.670,.745]} & +0.018 & 1.1 & 0.665\,{\scriptsize[.651,.678]} & 5.3 \\
GBT, all features (offline) & 0.729\,{\scriptsize[.717,.740]} & 0.675\,{\scriptsize[.643,.707]} & +0.007 & 3.9 & 0.674\,{\scriptsize[.661,.688]} & 12.2 \\
\midrule
Streaming GBT, never retrained & 0.604\,{\scriptsize[.587,.620]} & 0.594\,{\scriptsize[.557,.632]} & -0.117 & 13.0 & -- & -- \\
Streaming GBT, every 4 weeks & 0.717\,{\scriptsize[.702,.731]} & 0.714\,{\scriptsize[.675,.749]} & -0.005 & 1.2 & -- & -- \\
Streaming GBT, ADWIN & 0.729\,{\scriptsize[.712,.744]} & 0.705\,{\scriptsize[.665,.744]} & +0.008 & 1.3 & -- & -- \\
Streaming GBT, DDM & 0.724\,{\scriptsize[.708,.738]} & 0.694\,{\scriptsize[.656,.730]} & +0.003 & 1.3 & -- & -- \\
Streaming GBT, OPTWIN & 0.723\,{\scriptsize[.707,.737]} & 0.688\,{\scriptsize[.648,.723]} & +0.001 & 1.4 & -- & -- \\
Incremental trees, River ARF & 0.478\,{\scriptsize[.464,.492]} & 0.597\,{\scriptsize[.556,.641]} & -- & 28.0 & -- & -- \\
Incremental trees, River HAT & 0.332\,{\scriptsize[.322,.343]} & 0.393\,{\scriptsize[.352,.435]} & -- & 30.1 & -- & -- \\
Incremental trees, CapyMOA ARF & 0.527\,{\scriptsize[.518,.537]} & 0.610\,{\scriptsize[.579,.640]} & -- & 96.0 & -- & -- \\
\bottomrule
\end{tabular}

\end{table*}

\emph{Baselines.} The rule reaches 0.721 at 1\,h and 0.637 at 6\,h. A GBT given only total-power statistics does worse (0.695), a sign of the shift between the training and test years.

\emph{Offline model.} GBT on telemetry and peer deviations reaches 0.739 [0.726, 0.753] at 1\,h, better than the rule by +0.018 [+0.004, +0.032], and 0.665 at 6\,h (+0.028 [+0.007, +0.049]). Its clearest advantage is at the top of the ranking: to catch 25\% of onsets one hour ahead it raises 1.06 false alerts per node-day against 4.90 for the rule (paired difference +3.84 [+2.17, +4.22]), i.e.\ a precision of 1.2\%, 23 times the base rate.

\emph{Streaming.} Never retrained, the model is worse than the rule (0.604; 13 alerts). Refit every four weeks it reaches 0.717, statistically tied with the rule in AUC ($-$0.005 [$-$0.017, +0.007]) but with 1.15 alerts per node-day (+3.75 [+2.05, +4.12]). Triggering the refit with ADWIN, DDM or OPTWIN (0.729, 0.724, 0.723) gives the same AUC and 1.3--1.4 alerts, with 33, 52 and 58 retrains against 29 for the schedule (\cref{fig:retrains}): drift detection does not beat a fixed schedule here. Incremental trees fail (AUC 0.33--0.53).

\begin{figure}[t]
\centering
\begin{tikzpicture}
\begin{axis}[width=\columnwidth, height=3.6cm, xlabel={year}, ytick={1,2,3,4}, yticklabels={OPTWIN, DDM, ADWIN, every 4 weeks},
  ymin=0.4, ymax=4.6, xmin=2020.4, xmax=2022.8, tick label style={font=\scriptsize}, label style={font=\footnotesize},
  scaled x ticks=false, xticklabel style={/pgf/number format/.cd, fixed, precision=0}, xtick={2021,2022}]
\addplot[only marks, mark=|, mark size=3pt, thick, cmodel] table[col sep=comma, x=decyear, y=y] {results/retrains.csv};
\draw[dashed, crule] (axis cs:2022.23,0.4) -- (axis cs:2022.23,4.6) node[above left, font=\scriptsize, crule] {test};
\end{axis}
\end{tikzpicture}
\caption{Retraining times of the four policies. Drift detectors fire in bursts, mostly around maintenance events, and more often than the schedule without a gain.}
\label{fig:retrains}
\end{figure}

\emph{Node-specific onsets.} Removing cluster-wide onsets leaves 392 (1\,h) and 321 (6\,h) onsets. Here no learned model is better than the rule: the offline model scores 0.710 against 0.731 at 1\,h ($-$0.020 [$-$0.054, +0.016]) and 0.562 against 0.667 at 6\,h ($-$0.106 [$-$0.153, $-$0.058]) (\cref{tab:featuresets} lists the other feature sets). Calendar features alone fall below chance for node-specific onsets (0.373 at 6\,h), consistent with learning the maintenance schedule.

\begin{figure}[t]
\centering
\begin{tikzpicture}
\begin{axis}[width=\columnwidth, height=4.2cm, xlabel={recall of onsets}, ylabel={false alerts / node-day}, ymode=log,
  xmin=0, xmax=1, legend style={font=\scriptsize, at={(0.97,0.05)}, anchor=south east},
  tick label style={font=\scriptsize}, label style={font=\footnotesize}, every axis plot/.append style={thick}]
\addplot[cmodel, mark=none] table[col sep=comma, x=recall, y=model] {results/operating_curve.csv};
\addplot[crule, mark=none] table[col sep=comma, x=recall, y=rule] {results/operating_curve.csv};
\addplot[cgrey, dashed, mark=none] table[col sep=comma, x=recall, y=random] {results/operating_curve.csv};
\legend{GBT, rule, random}
\end{axis}
\end{tikzpicture}
\caption{False alerts per node-day needed to reach a given recall of 1\,h onsets. The rule's curve is flat at low recall because many idle windows tie at the same value.}
\label{fig:operating}
\end{figure}

\begin{table*}[t]
\caption{Offline GBT by feature set: mean per-node AUC for all and node-specific onsets.}
\label{tab:featuresets}
\centering\small
\begin{tabular}{lrrrr}
\toprule
Features & 1\,h all & 1\,h node-spec. & 6\,h all & 6\,h node-spec. \\
\midrule
Telemetry + peers & 0.739\,{\scriptsize[.726,.753]} & 0.710\,{\scriptsize[.670,.745]} & 0.665\,{\scriptsize[.651,.678]} & 0.562\,{\scriptsize[.524,.598]} \\
All features & 0.729\,{\scriptsize[.717,.740]} & 0.675\,{\scriptsize[.643,.707]} & 0.674\,{\scriptsize[.661,.688]} & 0.449\,{\scriptsize[.413,.486]} \\
Telemetry only & 0.717\,{\scriptsize[.703,.732]} & 0.711\,{\scriptsize[.669,.745]} & 0.607\,{\scriptsize[.590,.623]} & 0.585\,{\scriptsize[.549,.620]} \\
Calendar only & 0.537\,{\scriptsize[.529,.544]} & 0.502\,{\scriptsize[.472,.533]} & 0.565\,{\scriptsize[.552,.580]} & 0.373\,{\scriptsize[.331,.414]} \\
Rack context only & 0.508\,{\scriptsize[.498,.518]} & 0.489\,{\scriptsize[.467,.511]} & 0.497\,{\scriptsize[.491,.503]} & 0.511\,{\scriptsize[.492,.527]} \\
\bottomrule
\end{tabular}

\end{table*}

\emph{Operating point and stability.} \Cref{fig:operating} shows the cost of recall: half of the onsets cost about 14 alerts per node-day (precision 0.18\%), so the model is a risk ranking rather than an alarm. Results also depend strongly on which events fall in the test period (\cref{fig:monthly}): the pooled AUC of the offline model at 1\,h ranges from 0.54 to 0.98 across months, and months dominated by cluster-wide drains are not uniformly easy (0.98 and 0.93 in March and May, 0.64 and 0.54 in July and September).

\begin{figure}[t]
\centering
\begin{tikzpicture}
\begin{axis}[ybar, bar width=5pt, width=\columnwidth, height=4.2cm, ylabel={pooled AUC}, ymin=0.4, ymax=1.0,
  xtick=data, xticklabels from table={results/monthly_auc.csv}{label}, x tick label style={font=\scriptsize},
  tick label style={font=\scriptsize}, label style={font=\footnotesize},
  legend style={font=\scriptsize, at={(0.98,0.96)}, anchor=north east}]
\addplot[fill=cmodel, draw=none] table[col sep=comma, x expr=\coordindex, y=auc_tp_1h] {results/monthly_auc.csv};
\addplot[fill=crule, draw=none] table[col sep=comma, x expr=\coordindex, y=auc_tp_6h] {results/monthly_auc.csv};
\legend{1\,h, 6\,h}
\end{axis}
\end{tikzpicture}
\caption{Pooled AUC of the offline GBT per test month (2022). Asterisks mark months in which at least half of the onsets belong to cluster-wide events.}
\label{fig:monthly}
\end{figure}
